
\documentclass[11pt,a4paper]{article}
\usepackage[T1]{fontenc}
\usepackage{lmodern}
\usepackage{amsmath,amssymb,amsfonts,amsthm}
\usepackage{geometry}
\geometry{a4paper, margin=25mm}
\usepackage{booktabs}
\usepackage{xcolor}
\usepackage{fancyhdr}
\usepackage{hyperref}

\hypersetup{
    colorlinks=true,
    linkcolor=blue!60!black,
    citecolor=green!50!black,
    urlcolor=blue!80!black
}

\pagestyle{fancy}
\fancyhf{}
\rhead{\small CS/DT Coursework 2026}
\lhead{\small Deliverable 1: High-Level Programmatic Document}
\cfoot{\thepage}

\title{\textbf{Project Specification: Layer-Resolved Structural Aliasing and Cyber-Physical Vulnerabilities in Time-Series Foundation Models}}
\author{\textbf{Deliverable 1: High-Level Programmatic Document} \\
Joint Computer Security / Digital Transformation Graduate Program \\
Coursework Co-Supervision: Prof. Benini \\
Project Collaborators: Brignoli, Boniotti, and Sabbadini}
\date{\today}

\begin{document}

\maketitle

\section{Analyse Purpose and Audience (The `What?')}
This investigation formalizes the phenomenon of \textit{structural aliasing} within patch-based time-series foundation models, focusing specifically on Amazon's Chronos-Bolt framework. The physical boundary of this research is defined by the 1-minute multi-variate telemetry streams contained within \texttt{Dataset-SolarTechLab.csv}, prioritizing global tilted irradiance ($G_{\mathrm{tilt}}$) and active photovoltaic power generation ($P_{\mathrm{PV}}$). 

Unlike traditional macro-level evaluations of foundation model accuracy, this project implements a highly localized, layer-resolved evaluation blueprint across different stages of the Chronos backbone. By mapping representations across these specific layers, the project isolates where task-relevant signal attributes are preserved or compromised under architectural synchronization constraints. The investigation main goal is to evaluate the model using its standard baseline hyperparameter configuration, where patch size is set to $P=16$ tokens and stride is set to $S=16$ tokens.

\section{Audience Value and Target Impact (The `Why?')}
The intended audience for this research comprises automated smart-grid operations engineers, machine learning interpretability researchers, and academic evaluators led by Prof. Benini. The core objective is to challenge the dangerous assumption that pre-trained time-series large language models (LLMs) safely process sub-sampled continuous telemetry streams as long as the inputs satisfy the classical Nyquist-Shannon sampling theorem. 

We demonstrate that when a high-frequency solar event matches the model's fixed patch geometry, an architectural data deletion occurs that renders the physical event completely invisible to the self-attention mechanism. 

In automated micro-grid load-balancing settings, this representation collapse leads directly to semantic misclassifications, creating secondary vulnerabilities including automated battery over-discharge, flawed natural language explanations, and unmitigated cyber-physical stability risks. This documentation establishes the scene for our future research agenda, outlining known signal gaps and providing a robust conceptual foundation before empirical validation.

\section{Persuasion Strategy and Technical Architecture (The `How?')}
To construct a rigorous, falsifiable argument that challenges how practitioners trust time-series foundation models, we implement a multi-layered validation and defense pipeline:
\begin{itemize}
    \item \textbf{Hierarchical Spectral Partitioning}: We recursively divide the operational frequency band of solar transients into a multi-tiered hierarchy of binary classification bands to isolate decision boundary uncertainty and evaluate probe performance across fine-grained sub-bands.
    \item \textbf{Information-Theoretic Probing}: We deploy online Minimum Description Length (MDL) probes using continuous replay streams on the extracted hidden states of every individual decoder layer. We use the normalized Space Saving ($SV$) metric to quantify the exact linear decodability of signal attributes.
    \item \textbf{Causal Intervention}: We apply sequential Least-Squares Concept Erasure (LEACE) across the transformer stack to surgically ablate coarse binary frequency abstractions, measuring the exact causal impact on closed-loop autoregressive forecasting performance via Spectral Root Mean Squared Error (RMSE).
    \item \textbf{Joint Bayesian Priors}: We implement a probabilistic assessment engine using \texttt{PyMC}. We combine a highly objective \textit{Skeptical Prior} on the structural aliasing effect parameters ($\beta \sim \mathcal{N}(0, 0.1^2)$) with an \textit{Empirical Data-Calibrated Prior} representing real-world hardware instrument noise ($\sigma \sim \text{Half-Normal}(0.0616)$), calculated directly from the nighttime null intervals of \texttt{Dataset-SolarTechLab.csv}.
    \item \textbf{Dual-Engine Generative Testing}: We evaluate the model's manifold behavior using full \textit{KernelSynth} sum-product randomized kernel trees to model non-stationary background environments, while restricting the \textit{TSMixup} operator to run exclusively over pure sinusoids to maintain a clean algebraic control environment for phase-locking evaluation.
    \item \textbf{NIST AI RMF 1.0 Integration}: We frame the identified data-deletion mechanism as a formal zero-payload security vulnerability within the \textit{Map}, \textit{Measure}, and \textit{Manage} core functions of the NIST AI Risk Management Framework, evaluating whether proposed algorithmic defenses successfully reduce risk or leave significant residual exposure.
\end{itemize}

\section{Open Source and Licensing Compliance}
In strict compliance with the core open-source policies set forth by the coursework handbook, all project components are developed under a transparent public licensing framework. This technical programmatic document and the subsequent final report are released under the Creative Commons Attribution 4.0 International (CC-BY 4.0) license. All source code, probing scripts, and model intervention wrappers generated across the project lifetime are released under the open-source MIT license, ensuring full reproducibility, fairness of assessment, and academic integrity.

\end{document}
